\section{Discussion}
\label{sec:discussion}

\subsection{Lessons Learned}

Agent skills have become a security boundary between third party artifacts and
agent execution. A skill can combine natural language instructions, scripts,
dependencies, installation steps, update logic, and tool access within one
reusable package. Security relevant behavior can span several artifacts and
several stages of execution. Reviewing an isolated instruction, script, or
runtime event provides only part of the security context.

Our controlled and real world studies support a broader view of skill
security. Static analysis is effective for finding security relevant
capabilities at scale, while runtime evidence establishes which operations are
actually realized during execution. Policy context determines the security
meaning of the resulting behavior. Skill auditing benefits from preserving
these distinctions instead of collapsing capability, execution, and policy
into a single risk signal.

The real world audit also shows the importance of installation and maintenance
behavior. Confirmed malicious skills can place security relevant actions inside
workflows that appear operationally useful, including dependency acquisition,
updates, scheduled execution, and helper utilities. One confirmed sample
acquires content from a malicious external source, registers execution beyond
the initiating user interaction, and obtains authority to modify local agent
and skill state. The security meaning emerges from the complete path and its
relation to the function presented by the skill.

Skill ecosystems need review mechanisms that follow security relevant behavior
across the whole package and into execution. Marketplace screening can identify
candidates, while provenance and runtime evidence support deeper review of
behavior that crosses security boundaries. Clear evidence also makes findings
easier to verify, communicate, and act on by maintainers and platform
operators.


\subsection{Limitations}

\system observes bounded executions. A malicious path that depends on delayed
activation, unavailable external state, specific credentials, rare model
choices, or an untriggered instruction may remain outside the observed
execution. Coverage states record these execution boundaries and preserve the
path prefix recovered before stronger runtime support ends.

LLM mediated execution can also vary across replays. Different model decisions
or tool sequences may expose different portions of the same skill behavior.
Our evaluation uses fixed task inputs and model configurations for each
experiment, so the reported runtime provenance reflects behavior observed
under that execution setting. Repeated or diversity guided replay could extend
coverage for behavior that depends on execution choices.

Runtime observation is strongest for data carrying behavior that appears in
files, processes, tool calls, model context, and network activity. Local state
changes such as persistence updates, permission changes, and instruction state
modification are less consistently realized during bounded replay. Broader
execution environments and stronger triggering strategies can extend coverage
for these behaviors.

The real world audit covers 6,000 public skill samples collected from skills.sh
between January and August 2026. Other registries may differ in package
structure, moderation, permission models, runtime dependencies, and attack
distribution. The audit evaluates the operational use of \system for public
skill review and does not estimate malicious prevalence or recall for the
complete skill ecosystem.


\subsection{Responsible Disclosure}

Confirmed real world findings enter our disclosure process after manual review
of the original skill package and the evidence recovered by \system. The
review verifies the security relevant operation, supporting artifacts, runtime
behavior when available, external resources involved in the path, and the
resulting security effect before a finding is reported to a maintainer or
platform operator.

At the current snapshot, manual evidence review has confirmed 21 malicious
public skills from the dynamically analyzed candidates. Verification of
additional candidates is ongoing, and confirmed findings are being processed
through responsible disclosure channels. Disclosure records and platform
responses are maintained separately from automated predictions so that only
manually verified findings are reported externally.


\subsection{Ethical Considerations}

Our experiments are designed to study malicious behavior without exposing
users, credentials, or third party systems to the resulting effects. ProvBench
uses synthetic protected data and controlled runtime fixtures. Public skills
are executed inside an isolated analysis environment without real user
credentials. Security relevant files and secrets used during dynamic analysis
are synthetic, and external effects are constrained to controlled or simulated
resources where required.

The real world study analyzes publicly distributed skill artifacts. Collection
does not require interaction with skill authors or users, and the audit does
not intentionally exercise malicious actions against third party services.
Potentially harmful findings are manually verified before disclosure, and
operational details that could facilitate abuse are withheld while disclosure
is in progress. Our study follows established guidance for responsible
security and network measurement research
\cite{menlo,ethical}.